\chapter{Pendahuluan}

\section{Latar Belakang}

Pengelolaan aset dan inventori merupakan bagian penting dalam mendukung aktivitas operasional organisasi karena berkaitan dengan proses pencatatan, pemantauan, dan pengendalian sumber daya perusahaan \cite{ref3}. Pengelolaan aset yang baik dapat mendukung aktivitas operasional organisasi, khususnya dalam proses pemeliharaan, perbaikan, dan pengadaan aset \cite{ref4}. Untuk mendukung proses tersebut, organisasi mulai memanfaatkan sistem informasi manajemen aset yang mampu membantu proses pencatatan, pendataan, serta pengelolaan aset secara lebih terstruktur dan terdokumentasi \cite{ref4}.

Perkembangan teknologi informasi juga mendorong pertumbuhan \textit{platform} digital berbasis
\textit{Software as a Service} (SaaS) yang memungkinkan satu aplikasi digunakan oleh banyak
organisasi dalam satu layanan terpusat \cite{ref3}. Model layanan \textit{SaaS} memberikan
kemudahan bagi organisasi dalam menggunakan aplikasi tanpa harus membangun dan mengelola
infrastruktur sistem secara mandiri. Selain itu, kebutuhan terhadap sistem yang \textit{scalable}
dan mampu mendukung banyak organisasi menyebabkan pendekatan \textit{multitenancy} semakin banyak
diterapkan pada pengembangan aplikasi \textit{SaaS} modern \cite{ref4}.

Pendekatan \textit{multitenancy} memungkinkan satu aplikasi digunakan secara bersama oleh banyak
\textit{tenant} dengan tetap mempertahankan isolasi data masing-masing organisasi
\cite{ref3,ref5}. Dalam implementasinya, setiap \textit{tenant} dapat memiliki pengguna, data,
dan proses operasional yang berbeda-beda sehingga sistem membutuhkan arsitektur yang mampu
mengelola banyak \textit{tenant} tanpa mencampurkan data antar organisasi. Kondisi tersebut
menyebabkan aspek pengelolaan \textit{tenant} dan isolasi data menjadi bagian penting dalam
pengembangan aplikasi berbasis \textit{SaaS}.

Penelitian ini dilakukan pada pengembangan modul \textit{Asset Management}, \textit{Inventory},
dan \textit{Rental} berbasis \textit{web} yang terintegrasi dengan \textit{platform} RTPintar.
Modul yang dikembangkan dirancang untuk dapat digunakan oleh banyak organisasi dalam satu sistem
dengan kebutuhan pengelolaan data dan proses operasional yang berbeda pada setiap \textit{tenant}.

Dalam implementasi aplikasi \textit{SaaS}, terdapat beberapa pendekatan arsitektur \textit{tenant}
yang umum digunakan, seperti \textit{separate database}, \textit{separate schema}, maupun
\textit{shared schema} \cite{ref6}. Pendekatan \textit{separate database} menawarkan tingkat
isolasi data yang tinggi karena setiap \textit{tenant} menggunakan basis data yang terpisah.
Namun, pendekatan tersebut menyebabkan setiap \textit{tenant} memerlukan pengelolaan dan
konfigurasi basis data tersendiri sehingga meningkatkan kompleksitas operasional ketika jumlah
\textit{tenant} terus bertambah \cite{ref7}. Selain itu, pendekatan \textit{database-per-tenant}
juga meningkatkan kebutuhan \textit{resource} infrastruktur dan proses \textit{maintenance} karena
setiap \textit{tenant} memiliki lingkungan penyimpanan data yang terpisah \cite{ref7}.

Sebagai solusi terhadap permasalahan tersebut, arsitektur \textit{multitenant} dengan pendekatan
\textit{shared database} dan \textit{shared schema} memungkinkan seluruh \textit{tenant}
menggunakan basis data dan struktur tabel yang sama \cite{ref6}. Dengan pendekatan tersebut,
\textit{tenant} baru dapat ditambahkan tanpa perlu membuat basis data baru untuk setiap
\textit{tenant} sehingga proses \textit{provisioning} dapat dilakukan lebih cepat \cite{ref8}.
Selain itu, penggunaan \textit{shared database} dan \textit{shared schema} lebih sesuai untuk
aplikasi \textit{SaaS} dengan jumlah \textit{tenant} yang dapat terus bertambah karena pengelolaan
\textit{tenant} dan \textit{resource} sistem dapat dilakukan secara lebih terpusat \cite{ref7}.
Oleh karena itu, pendekatan \textit{shared database} dan \textit{shared schema} dipilih pada
penelitian ini untuk mendukung kebutuhan pengelolaan banyak \textit{tenant} dalam satu
\textit{platform}.

Meskipun memiliki kelebihan dalam penggunaan \textit{resource} yang lebih efektif, implementasi
\textit{shared database} dan \textit{shared schema} memiliki tantangan pada aspek isolasi data
antar \textit{tenant} \cite{ref9}. Penggunaan struktur tabel yang sama menyebabkan seluruh data
\textit{tenant} berada dalam basis data yang sama sehingga kesalahan dalam pengelolaan
\textit{tenant context} maupun kontrol akses dapat menyebabkan terjadinya akses lintas
\textit{tenant} yang tidak sesuai \cite{ref6}. Kondisi tersebut berpotensi menimbulkan kebocoran
data dan ketidaksesuaian akses informasi antar \textit{tenant} apabila mekanisme isolasi
\textit{tenant} tidak diterapkan secara konsisten pada lapisan aplikasi maupun basis data
\cite{ref4} \cite{ref9}.

Berdasarkan kondisi tersebut, penelitian ini berfokus pada perancangan arsitektur
\textit{multitenant} berbasis \textit{shared database} dan \textit{shared schema} pada modul
\textit{Asset Management}, \textit{Inventory}, dan \textit{Rental} berbasis \textit{web} yang
terintegrasi dengan \textit{platform} RTPintar. Penelitian mencakup perancangan mekanisme
\textit{tenant identification} berbasis \textit{middleware}, implementasi \textit{tenant scoping}
pada lapisan aplikasi, perancangan model penyimpanan data \textit{multitenant}, serta proses
\textit{provisioning tenant} otomatis untuk mendukung proses \textit{onboarding tenant} baru.

\section{Rumusan Masalah}

Berdasarkan latar belakang yang telah dipaparkan, maka rumusan masalah dalam penelitian ini
adalah sebagai berikut:

\begin{enumerate}
  \item Bagaimana merancang arsitektur \textit{multitenant} berbasis \textit{shared database} dan
    \textit{shared schema} yang mampu meningkatkan efisiensi infrastruktur serta mendukung
    pengelolaan banyak organisasi dalam satu \textit{platform}?
  \item Bagaimana merancang mekanisme isolasi data dan \textit{provisioning} \textit{tenant} otomatis
    pada arsitektur \textit{multitenant} agar proses \textit{onboarding} organisasi baru dapat dilakukan
    secara konsisten, aman, dan tanpa konfigurasi manual tambahan?
\end{enumerate}

\section{Tujuan dan Manfaat}

Penelitian ini bertujuan untuk merancang arsitektur \textit{multitenant} berbasis \textit{shared database}
dan \textit{shared schema} yang mampu mendukung pengelolaan banyak organisasi dalam satu \textit{platform}
secara efisien. Perancangan arsitektur dilakukan untuk mengatasi keterbatasan pendekatan non
\textit{multitenant} yang menyebabkan kompleksitas pengelolaan sistem dan peningkatan kebutuhan
infrastruktur ketika jumlah \textit{tenant} bertambah. Selain itu, penelitian ini juga bertujuan merancang
mekanisme isolasi data antar \textit{tenant} melalui \textit{tenant scoping} pada lapisan aplikasi serta
merancang proses \textit{provisioning} \textit{tenant} otomatis agar \textit{onboarding} organisasi baru
dapat dilakukan secara konsisten, aman, dan tanpa konfigurasi manual tambahan.

Melalui penelitian ini diharapkan rancangan arsitektur yang dihasilkan dapat memberikan manfaat
praktis dalam mendukung pengembangan \textit{platform} \textit{multi organization} berbasis \textit{web}, khususnya
pada sistem manajemen aset, inventori, dan rental. Pendekatan \textit{multitenant} berbasis
\textit{shared database} dan \textit{shared schema} diharapkan mampu meningkatkan efisiensi
penggunaan \textit{resource} infrastruktur, menyederhanakan pengelolaan \textit{tenant}, serta mendukung
proses \textit{onboarding} \textit{tenant} secara lebih cepat dan terstandarisasi. Selain manfaat praktis,
penelitian ini juga diharapkan dapat menjadi referensi akademik dalam pengembangan penelitian
terkait arsitektur \textit{multitenant}, \textit{tenant isolation}, dan \textit{provisioning} \textit{tenant} pada
sistem \textit{Software as a Service} (\textit{SaaS}).

Keterkaitan antara tujuan penelitian, pengujian yang dilakukan, dan kesimpulan yang diharapkan
disajikan pada Tabel \ref{tab:tujuan_pengujian}.

\begin{table}[H]
  \caption{Tabel keterkaitan antara tujuan, pengujian dan kesimpulan.}
  \label{tab:tujuan_pengujian}
  \centering
  \fontsize{10}{12}\selectfont
  \setlength{\arrayrulewidth}{0.5pt}
  \setlength{\tabcolsep}{4pt}
  \begin{tabular}{|>{\centering\arraybackslash}p{0.5cm}|>{\raggedright\arraybackslash}p{4.3cm}|>{\raggedright\arraybackslash}p{3.5cm}|>{\raggedright\arraybackslash}p{4.3cm}|}
    \hline
    \multicolumn{1}{|>{\centering\arraybackslash}m{0.5cm}|}{\cellcolor{black}\color{white}\textbf{No.}} &
    \multicolumn{1}{>{\centering\arraybackslash}m{4.3cm}|}{\cellcolor{black}\color{white}\textbf{Tujuan}} &
    \multicolumn{1}{>{\centering\arraybackslash}m{3.5cm}|}{\cellcolor{black}\color{white}\textbf{Pengujian}} &
    \multicolumn{1}{>{\centering\arraybackslash}m{4.3cm}|}{\cellcolor{black}\color{white}\textbf{Kesimpulan}} \\
    \hline
    1 &
    Merancang arsitektur \textit{multitenant} berbasis \textit{shared database} dan \textit{shared schema} untuk mendukung pengelolaan banyak organisasi dalam satu \textit{platform} secara efisien &
    Evaluasi efisiensi infrastruktur dan struktur arsitektur \textit{multitenant} &
    Sistem mampu menggunakan satu basis data dan satu skema bersama untuk banyak \textit{tenant} tanpa memerlukan \textit{instance} terpisah \\
    \hline
    2 &
    Merancang mekanisme isolasi data antar \textit{tenant} pada arsitektur \textit{multitenant} &
    Pengujian \textit{tenant isolation} dan \textit{tenant scoping} pada lapisan aplikasi &
    Data antar \textit{tenant} tetap terisolasi dan \textit{tenant} tidak dapat mengakses data milik \textit{tenant} lain \\
    \hline
    3 &
    Merancang mekanisme \textit{provisioning} \textit{tenant} otomatis untuk mendukung \textit{onboarding} \textit{tenant} baru secara lebih efisien &
    Pengujian proses \textit{provisioning} \textit{tenant} otomatis &
    \textit{Tenant} baru dapat dibuat beserta konfigurasi awalnya tanpa proses konfigurasi manual tambahan \\
    \hline
    4 &
    Merancang arsitektur yang mendukung konsistensi pengelolaan sistem \textit{multi organization} &
    Evaluasi konsistensi struktur data dan pengelolaan \textit{tenant} &
    Setiap \textit{tenant} memiliki struktur pengelolaan data dan konfigurasi yang konsisten dalam satu \textit{platform} terpusat \\
    \hline
  \end{tabular}
\end{table}

\section{Batasan Masalah}

Penelitian ini difokuskan pada perancangan arsitektur \textit{multitenant} berbasis \textit{shared database}
dan \textit{shared schema} pada \textit{platform} manajemen aset, inventori, dan rental berbasis \textit{web}. Agar
penelitian lebih terarah dan sesuai dengan ruang lingkup tugas akhir, maka ditetapkan beberapa
batasan masalah sebagai berikut:

\begin{enumerate}
  \item Penelitian difokuskan pada perancangan arsitektur sistem \textit{multitenant} dan tidak membahas
    implementasi penuh seluruh fitur aplikasi pada \textit{platform} manajemen aset, inventori, dan rental.
  \item Model \textit{multitenancy} yang digunakan dibatasi pada pendekatan \textit{shared
    database} dan \textit{shared schema} tanpa membandingkan implementasi secara mendalam dengan
    model \textit{separate database} atau \textit{separated schema}.
  \item Mekanisme isolasi data \textit{tenant} difokuskan pada penerapan \textit{tenant scoping} pada
    lapisan aplikasi dan penggunaan \textit{tenant identifier} pada basis data.
  \item Penelitian tidak membahas mekanisme keamanan lanjutan seperti enkripsi tingkat lanjut,
    \textit{distributed security architecture}, maupun audit keamanan infrastruktur secara
    menyeluruh.
  \item Proses \textit{provisioning} \textit{tenant} dibatasi pada pembuatan \textit{tenant} baru, inisialisasi
    konfigurasi awal \textit{tenant}, serta pengelolaan data dasar yang diperlukan dalam sistem.
  \item Evaluasi penelitian difokuskan pada aspek efisiensi infrastruktur, konsistensi isolasi
    data antar \textit{tenant}, dan efektivitas \textit{provisioning} \textit{tenant} otomatis.
  \item Penelitian difokuskan pada konteks \textit{platform} \textit{multi organization} berbasis \textit{web} dan
    tidak membahas implementasi pada arsitektur \textit{distributed system} atau
    \textit{multi region infrastructure}.
  \item Rancangan arsitektur yang dihasilkan ditujukan sebagai dasar implementasi pada proyek
    \textit{capstone} dan tidak membahas \textit{deployment} sistem pada lingkungan produksi
    berskala besar secara penuh.
\end{enumerate}

\section{Metode Penelitian}

Pendekatan yang digunakan dalam penelitian ini mencakup studi literatur, perancangan sistem,
implementasi, serta pengujian sistem. Pendekatan tersebut digunakan untuk mendukung proses
perancangan arsitektur \textit{multitenant} berbasis \textit{shared database} dan \textit{shared schema}
pada modul \textit{Asset Management}, \textit{Inventory}, dan \textit{Rental} berbasis \textit{web}.

Tahapan penelitian dilakukan melalui proses analisis kebutuhan sistem, perancangan arsitektur
\textit{multitenant}, implementasi mekanisme \textit{tenant isolation} dan \textit{provisioning} \textit{tenant},
serta pengujian sistem untuk memastikan pengelolaan data antar \textit{tenant} berjalan sesuai dengan
rancangan yang telah dibuat.

\section{Jadwal Pelaksanaan}

Jadwal pelaksanaan penelitian disusun berdasarkan tahapan perancangan arsitektur \textit{multitenant} yang
meliputi studi literatur, analisis kebutuhan sistem dan identifikasi permasalahan arsitektural
pada pendekatan non \textit{multitenant}, perancangan arsitektur \textit{multitenant} berbasis \textit{shared
database} dan \textit{shared schema}, perancangan mekanisme isolasi data dan \textit{provisioning}
\textit{tenant} otomatis, implementasi prototipe arsitektur sebagai pembuktian konsep, evaluasi kesesuaian
arsitektur terhadap parameter yang telah ditetapkan, serta penyusunan laporan akhir penelitian.
Rincian waktu dan tahapan pelaksanaan penelitian disajikan pada Tabel \ref{tab:jadwal}.

\begin{table}[H]
  \caption{Jadwal Pelaksanaan Tugas Akhir.}
  \label{tab:jadwal}
  \centering
  \fontsize{10}{12}\selectfont
  \setlength{\arrayrulewidth}{0.5pt}
  \setlength{\tabcolsep}{4pt}
  \begin{tabular}{|>{\centering\arraybackslash}p{0.4cm}|>{\raggedright\arraybackslash}p{4.8cm}|>{\centering\arraybackslash}p{1.0cm}|>{\centering\arraybackslash}p{1.0cm}|>{\centering\arraybackslash}p{1.0cm}|>{\centering\arraybackslash}p{1.0cm}|>{\centering\arraybackslash}p{1.0cm}|>{\centering\arraybackslash}p{1.0cm}|}
    \hline
    \multicolumn{1}{|>{\centering\arraybackslash}m{0.4cm}|}{\cellcolor{black}\color{white}\textbf{No.}} &
    \multicolumn{1}{>{\centering\arraybackslash}m{4.8cm}|}{\cellcolor{black}\color{white}\textbf{Deskripsi Tahapan}} &
    \multicolumn{1}{>{\centering\arraybackslash}m{1.0cm}|}{\cellcolor{black}\color{white}\textbf{\shortstack{Bulan\\1}}} &
    \multicolumn{1}{>{\centering\arraybackslash}m{1.0cm}|}{\cellcolor{black}\color{white}\textbf{\shortstack{Bulan\\2}}} &
    \multicolumn{1}{>{\centering\arraybackslash}m{1.0cm}|}{\cellcolor{black}\color{white}\textbf{\shortstack{Bulan\\3}}} &
    \multicolumn{1}{>{\centering\arraybackslash}m{1.0cm}|}{\cellcolor{black}\color{white}\textbf{\shortstack{Bulan\\4}}} &
    \multicolumn{1}{>{\centering\arraybackslash}m{1.0cm}|}{\cellcolor{black}\color{white}\textbf{\shortstack{Bulan\\5}}} &
    \multicolumn{1}{>{\centering\arraybackslash}m{1.0cm}|}{\cellcolor{black}\color{white}\textbf{\shortstack{Bulan\\6}}} \\
    \hline
    1 & Studi Literatur dan Pengumpulan Referensi & \cellcolor{black} & & & & & \tabularnewline
    \hline
    2 & Analisis Kebutuhan Sistem dan Identifikasi Permasalahan & \cellcolor{black} & \cellcolor{black} & & & & \tabularnewline
    \hline
    3 & Perancangan Arsitektur Multitenant & & \cellcolor{black} & \cellcolor{black} & & & \tabularnewline
    \hline
    4 & Perancangan Mekanisme \textit{Tenant Isolation} dan \textit{Provisioning Tenant} & & & \cellcolor{black} & \cellcolor{black} & & \tabularnewline
    \hline
    5 & Implementasi Arsitektur dan Fitur \textit{Multitenancy} & & & & \cellcolor{black} & \cellcolor{black} & \tabularnewline
    \hline
    6 & Evaluasi dan Pengujian Sistem & & & & & \cellcolor{black} & \tabularnewline
    \hline
    7 & Analisis Hasil Pengujian & & & & & \cellcolor{black} & \cellcolor{black} \tabularnewline
    \hline
    8 & Penyusunan dan Revisi Laporan Tugas Akhir & & & & & & \cellcolor{black} \tabularnewline
    \hline
    9 & Persiapan Presentasi dan Sidang Tugas Akhir & & & & & & \cellcolor{black} \tabularnewline
    \hline
  \end{tabular}
\end{table}
